\documentclass[12pt]{article}
\usepackage[utf8]{inputenc}
\usepackage[T1]{fontenc}
\usepackage{hyperref}
\title{Yaweyuha Language Resources for K-12 Curriculum Development}
\author{Compiled by Claude Sonnet 5.5 (claude-sonnet-5-5) for Andr\'as Kornai}
\date{}
\begin{document}
\maketitle

\section{Introduction}
Yaweyuha (also Yabiyufa or Yawiyuha; ISO 639-3: yby) is spoken in the Eastern Highlands of Papua New Guinea. Sources place it in the Kainantu-Goroka branch of Trans-New Guinea. Speaker counts conflict. Wikipedia's infobox gives 2,000 (1991), and its body text gives around 4,600 with a ``citation needed'' tag.

\section{Primary Educational Resources}
The search found no primers, textbooks or school programmes for Yaweyuha.

\section{Linguistic and Reference Materials}
\begin{itemize}
\item Wikipedia overview: \url{https://en.wikipedia.org/wiki/Yaweyuha_language}
\item Numeral system notes: \url{https://lingweb.eva.mpg.de/channumerals/Yaweyuha.htm}
\item Wikitongues language page (gives the writing system as Latin script): \url{https://wikitongues.org/languages/yby/}
\end{itemize}

\section{Religious and Community Resources}
A Yaweyuha New Testament exists. It was produced by Wycliffe Bible Translators, with copyright held by Biblica (1983, 2011). eBible.org lists it under a Creative Commons licence (BY-NC-ND 4.0).
\begin{itemize}
\item eBible.org record: \url{https://ebible.org/details.php?id=yby}
\item About page for the New Testament: \url{https://ebible.org/study/content/texts/yby/about.html}
\item PNG Bible record: \url{https://png.bible/details.php?id=yby}
\end{itemize}

\section{Digital and Audio Resources}
\begin{itemize}
\item Global Recordings Network: \url{https://globalrecordings.net/en/language/yby}
\end{itemize}
I did not check which recordings are available on that page.

\section{Curriculum Development Context}
The New Testament is the largest known text in Yaweyuha and shows the language has a working orthography. It could serve as a reading corpus. No primers, grammar or dictionary turned up in this search, so a curriculum would need new materials. Numeracy materials would need to consider whether speakers use traditional numerals or Tok Pisin ones beyond five.

\end{document}